\begin{table}[H]
\centering
\caption{Comparison of LSQ returns with S\&P 500 returns}
\label{tab:updated-table-2}
\begin{tabular}{lcc}
\toprule
\textbf{Metric} & \textbf{LSQ} & \textbf{S\&P 500 Index} \\
\midrule
\multicolumn{3}{l}{\emph{Statistical properties of monthly returns}} \\
\midrule
Average (arithmetic) return & 3.52 & 1.30 \\
Average (geometric) return & 3.45 & 1.19 \\
Percent Positive Months & 95.65 & 69.57 \\
Percent Negative Months & 4.35 & 30.43 \\
Average Positive Months & 3.72 & 3.62 \\
Average Negative Months & -0.87 & -4.01 \\
Best Month & 21.97 & 12.82 \\
Worst Month & -2.20 & -12.35 \\
\midrule
\multicolumn{3}{l}{\emph{First four moments of distribution of monthly returns}} \\
\midrule
Mean & 3.52 & 1.30 \\
Standard deviation & 3.98 & 4.51 \\
Skewness & 2.64 & -0.43 \\
Excess kurtosis & 8.54 & 0.54 \\
\midrule
\multicolumn{3}{l}{\emph{Risk adjusted performance metrics}} \\
\midrule
Sharpe Ratio & 0.88 & 0.29 \\
Sortino Ratio & 14.58 & 0.47 \\
Omega Ratio & 94.26 & 2.06 \\
Alpha & 3.55 & 0 \\
Beta & -0.03 & 1 \\
\bottomrule
\end{tabular}

\vspace{0.2cm}
\begin{minipage}{0.92\linewidth}
\footnotesize
Note: Returns are monthly percentages except ratios, skewness, excess kurtosis, and beta.
The LSQ column uses the baseline net-of-fees series.
The S\&P 500 benchmark is the total return index, consistent with the paper's dividend-reinvestment convention.
\end{minipage}
\end{table}
